%% ============================================================
%%  Conclusion et Perspectives
%% ============================================================
\chapter*{Conclusion et Perspectives}
\addcontentsline{toc}{chapter}{Conclusion et Perspectives}
\markboth{Conclusion et Perspectives}{Conclusion et Perspectives}

\section*{Rappel de la problématique}

Ce projet de fin d'études, réalisé au sein de la Salle des Marchés de BMCE Capital,
a été structuré autour de deux problèmes complémentaires~:

\begin{enumerate}
    \item \textbf{Comment identifier les meilleures opportunités de marché} parmi les
    93 titres du MASI avec suffisamment de rapidité et de traçabilité pour être
    actionnable en conditions réelles de trading~?
    \item \textbf{Comment valider rigoureusement des stratégies de trading} sans
    tomber dans le piège de l'overfitting, du data-snooping ou du biais de sélection
    sur tests multiples~?
\end{enumerate}

\section*{Synthèse des contributions}

\textbf{Contribution 1 — Signal Engine (pipeline A→G).}
Un moteur de signaux en sept couches transforme jusqu'à 30 candidats par famille d'indicateurs par horizon
en un score d'ensemble pondéré par robustesse temporelle. Chaque couche intègre un mécanisme
anti-overfitting distinct~: évaluation hors-échantillon systématique, gate de viabilité,
score multi-composante (Sharpe, stabilité, cohérence, drawdown), floor absolu,
clustering de corrélation de Pearson et moyenne pondérée.

\textbf{Contribution 2 — Moteur de Backtest WFO.}
Une implémentation fidèle de la Walk-Forward Analysis de Pardo (2008) avec la fonction
objectif PROM, une validation statistique par Deflated Sharpe Ratio et simulation
Monte Carlo, et un sizing calibré sur les trades OOS. Un kernel Numba JIT atteint un
speedup de $38\times$, rendant l'optimisation de l'ensemble du MASI viable en
conditions opérationnelles.

\textbf{Contribution 3 — Tableau de bord décisionnel.}
Un dashboard unifié qui consomme les deux moteurs pour répondre aux trois axes du
besoin desk~: rapidité (Score Opportunité en un coup d'œil), traçabilité (Section B
— Justification avec drill-down) et action (Section C — Niveaux stop/TP + narratif
conditionnel).

\textbf{Contribution 4 — Architecture industrialisable.}
Une plateforme full-stack (Next.js, FastAPI, PostgreSQL, MinIO, RQ) organisée en
quatre stations conformes au paradigme méta-stratégie de López de Prado, avec
87+ tests automatisés et traçabilité complète des runs.

\section*{Limites du travail}

L'honnêteté scientifique impose d'identifier les limites de l'approche~:

\begin{itemize}
    \item \textbf{Deflated Sharpe Ratio et CPCV}~: bien que le DSR soit calculé
    en post-validation, la Combinatorially Purged Cross-Validation (CPCV) de López
    de Prado — qui élimine la fuite de données temporelle dans la cross-validation —
    n'est pas encore intégrée au Signal Engine.
    \item \textbf{Fréquence journalière uniquement}~: le pipeline est conçu pour
    les données \textit{1D}. Les opportunités intraday et multi-fréquences ne sont
    pas couvertes.
    \item \textbf{Univers MASI uniquement}~: les produits dérivés (futures, options),
    le marché obligataire et les devises sont hors périmètre.
    \item \textbf{Portfolio oversight absent}~: la 5ème station de López de Prado
    (allocation dynamique, monitoring de drift, décommissionnement de stratégies)
    n'est pas implémentée.
    \item \textbf{Régimes non conditionnels}~: les signaux sont calculés sans
    condition explicite sur le régime de marché (trend / mean-revert / vol élevée).
\end{itemize}

\section*{Perspectives d'amélioration}

\textbf{Hypothesis Lab.}
Implémenter un classifieur de régime de marché (trend / mean-revert / volatilité)
permettant le \emph{déploiement conditionnel} des stratégies~: une stratégie SMA
n'est activée que si le régime est tendanciel. Semaine~5~\cite{rapport_s5} a esquissé
ce cadre~; il reste à le formaliser et l'intégrer.

\textbf{Feature Importance Ranking.}
López de Prado (2020, Ch.~6)~\cite{lopezdeprado2020} propose des méthodes basées
sur les forêts aléatoires (Mean Decrease Impurity, Mean Decrease Accuracy) pour
identifier quels indicateurs contribuent réellement au pouvoir prédictif à l'échelle
du portefeuille.

\textbf{Extension multi-instruments.}
Le pricer de futures développé en Semaines~8--9~\cite{rapport_s9} pose la base d'une
extension vers les produits dérivés marocains. Les données obligataires et devises
constituent des sources complémentaires pour enrichir le contexte macroéconomique.

\textbf{Portfolio Oversight.}
Allocation dynamique par Hierarchical Risk Parity (HRP) ou Markowitz~\cite{markowitz1952},
monitoring en temps réel du drift des métriques OOS, décommissionnement automatique
des stratégies dont le WFE se dégrade.

\textbf{Fréquences intraday.}
Extension du pipeline à des données 1H ou 15~min permettrait de couvrir les
opportunités de trading intrajournalier, particulièrement pertinentes pour les
périodes de publication de résultats ou d'annonces macroéconomiques.

\medskip

\noindent En définitive, ce projet démontre qu'il est possible de construire, dans un
délai de stage, une plateforme quantitative de niveau \og research-grade \fg{} sur le
marché marocain — un marché sous-couvert par les outils institutionnels existants.
Le Signal Engine et le moteur WFO constituent une fondation méthodologique solide,
prête à être enrichie et déployée progressivement à la Salle des Marchés de BMCE Capital.
